% Định nghĩa màu sắc nhận diện sách Stewart Calculus (nếu chưa được khai báo ở preamble)
\providecolor{stewartcyan}{RGB}{0, 130, 195}
\providecolor{stewartpurple}{RGB}{85, 37, 131}
\providecolor{stewartlightblue}{RGB}{232, 238, 245}
\providecolor{stewartred}{RGB}{204, 34, 41}

\fancyhead[L]{}
\fancyhead[R]{}
\fancyfoot[L]{\textbf{\textsf{74}}}
\fancyfoot[C]{}
\fancyfoot[R]{}

\noindent
\begin{minipage}[t]{0.26\textwidth}
    \vspace{0.2cm}
    % Ghi chú bên lề 1
    \noindent
    \colorbox{stewartlightblue}{%
        \parbox{\dimexpr\linewidth-2\fboxsep}{%
            \raggedright
            \colorbox{stewartpurple}{\textcolor{white}{\textbf{\textsf{\footnotesize PS}}}}%
            \hspace{0.25em}%
            \textsf{\footnotesize\textbf{Tương tự:} Thử với một bài toán tương tự, đơn giản hơn.}%
        }%
    }

    \vspace{5.9cm}
    % Ghi chú bên lề 2
    \noindent
    \colorbox{stewartlightblue}{%
        \parbox{\dimexpr\linewidth-2\fboxsep}{%
            \raggedright
            \colorbox{stewartpurple}{\textcolor{white}{\textbf{\textsf{\footnotesize PS}}}}%
            \hspace{0.25em}%
            \textsf{\footnotesize Tìm kiếm một quy luật.}%
        }%
    }
\end{minipage}%
\hfill
\begin{minipage}[t]{0.71\textwidth}
    \noindent
    \textbf{\textsf{\color{stewartcyan}VÍ DỤ 3}}\quad Nếu $f_0(x) = x/(x + 1)$ và $f_{n+1} = f_0 \circ f_n$ với $n = 0, 1, 2, \dots$, hãy tìm một công thức cho $f_n(x)$.

    \medskip
    \noindent
    \textbf{\textsf{\color{stewartcyan}LỜI GIẢI}}\quad Trước hết, ta tìm công thức của $f_n(x)$ cho các trường hợp riêng $n = 1, 2$ và $3$.
    \[\begin{aligned}
        f_1(x) &= (f_0 \circ f_0)(x) = f_0(f_0(x)) = f_0\left(\frac{x}{x + 1}\right) \\[6pt]
               &= \frac{\dfrac{x}{x + 1}}{\dfrac{x}{x + 1} + 1} = \frac{\dfrac{x}{x + 1}}{\dfrac{2x + 1}{x + 1}} = \frac{x}{2x + 1} \\[10pt]
        f_2(x) &= (f_0 \circ f_1)(x) = f_0(f_1(x)) = f_0\left(\frac{x}{2x + 1}\right) \\[6pt]
               &= \frac{\dfrac{x}{2x + 1}}{\dfrac{x}{2x + 1} + 1} = \frac{\dfrac{x}{2x + 1}}{\dfrac{3x + 1}{2x + 1}} = \frac{x}{3x + 1} \\[10pt]
        f_3(x) &= (f_0 \circ f_2)(x) = f_0(f_2(x)) = f_0\left(\frac{x}{3x + 1}\right) \\[6pt]
               &= \frac{\dfrac{x}{3x + 1}}{\dfrac{x}{3x + 1} + 1} = \frac{\dfrac{x}{3x + 1}}{\dfrac{4x + 1}{3x + 1}} = \frac{x}{4x + 1}
    \end{aligned}\]

    \noindent
    Ta nhận thấy một quy luật: hệ số của $x$ trong mẫu thức của $f_n(x)$ là $n + 1$ trong cả ba trường hợp ta vừa tính toán. Vì vậy, ta dự đoán rằng, một cách tổng quát,

    \medskip
    \noindent
    \parbox{\linewidth}{%
        \makebox[0pt][l]{\setlength{\fboxrule}{1pt}\fcolorbox{stewartred}{white}{\bfseries\sffamily\color{stewartred}\footnotesize 4}}%
        \centering
        $\displaystyle f_n(x) = \frac{x}{(n + 1)x + 1}$
    }

    \medskip
    \noindent
    Để chứng minh điều này, ta sử dụng Nguyên lý Quy nạp Toán học. Ta đã kiểm chứng rằng (4) đúng với $n = 1$. Giả sử rằng hệ thức này đúng với $n = k$, nghĩa là
    \[
        f_k(x) = \frac{x}{(k + 1)x + 1}
    \]
    Khi đó
    \[\begin{aligned}
        f_{k+1}(x) &= (f_0 \circ f_k)(x) = f_0(f_k(x)) = f_0\left(\frac{x}{(k + 1)x + 1}\right) \\[6pt]
                   &= \frac{\dfrac{x}{(k + 1)x + 1}}{\dfrac{x}{(k + 1)x + 1} + 1} = \frac{\dfrac{x}{(k + 1)x + 1}}{\dfrac{(k + 2)x + 1}{(k + 1)x + 1}} = \frac{x}{(k + 2)x + 1}
    \end{aligned}\]

    \noindent
    Biểu thức này cho thấy (4) đúng với $n = k + 1$. Do đó, theo nguyên lý quy nạp toán học, công thức đúng với mọi số nguyên dương $n$. \hfill{\color{stewartcyan}$\blacksquare$}
\end{minipage}